\noindent
Let $S$ be a scheme.
Let $(U, R, s, t, c)$ be a groupoid scheme over $S$.
Assume $U = \Spec(A)$, and $R = \Spec(B)$ are affine.
In this case we get two ring maps
$s^\sharp, t^\sharp : A \longrightarrow B$.
Let $C$ be the equalizer of $s^\sharp$ and $t^\sharp$. In a formula
\begin{equation}
\label{equation-invariants}
C = \{a \in A \mid t^\sharp(a) = s^\sharp(a) \}.
\end{equation}
We will sometimes call this the {\it ring of $R$-invariant functions on $U$}.
What properties does $M = \Spec(C)$ have? The first observation is
that the diagram
$$
\xymatrix{
R \ar[r]_s \ar[d]_t & U \ar[d] \\
U \ar[r] & M
}
$$
is commutative, i.e., the morphism $U \to M$ equalizes $s, t$.
Moreover, if $T$ is any affine scheme, and if $U \to T$ is
a morphism which equalizes $s, t$, then $U \to T$ factors through $U \to M$.
In other words, $U \to M$ is a coequalizer in the category of affine schemes.